%=====================================================================
\section{Final --- 17/07/2023}
%=====================================================================

\subsection{Ejercicio 1 --- Dispositivos de medición de agua (AES-ECB, HTTPS, integridad)}

\begin{consigna}
Hay una empresa que se encarga de controlar el agua. Para esto la empresa implementa unos
dispositivos encargados de hacer las respectivas mediciones. Los mismos le informan a la
central los datos recolectados.

El procedimiento cuando se pone un dispositivo es el siguiente: primero el dispositivo emite
su ID por broadcast, luego la central responde a ese ID y le envía una llave por HTTP. Cuando
el dispositivo le debe enviar lo recolectado a la empresa usa AES-ECB como método de
encripción.

La empresa contrata una auditoría y le dice que posee fallas en la comunicación de la llave y
problemas de integridad en la comunicación de los dispositivos.
\begin{enumerate}
  \item[a)] Usando HTTPS \textquestiondown{}se hubieran solucionado todos los problemas que informó la
  auditoría?
  \item[b)] En caso de no poder usar HTTPS por cuestiones de presupuestos \textquestiondown{}qué otro protocolo
  utilizaría para el intercambio de claves?
  \item[c)] \textquestiondown{}Qué cambios son necesarios para asegurar la integridad en el envío de información
  de los dispositivos hacia la central?
\end{enumerate}
\end{consigna}

\paragraph{Temas.}
Primera Parte, Clase 2 --- Cifrado (\S ECB, \S modos de cifrado, \S CPA-security); Clase 3 ---
MACs y Cifrado Autenticado (\S qué es un MAC, \S cifrado autenticado, \S modos AEAD:
CCM/GCM); Clase 5 --- Protocolos (\S TLS Handshake, \S Diffie--Hellman); Segunda Parte,
Clase 11 --- Protección de Datos (\S protocolos de transporte cifrados: TLS/HTTPS).

\paragraph{Qué necesitás saber.}
\begin{itemize}
  \item \textbf{ECB (Electronic Codebook):} cifra cada bloque de forma independiente,
  $c_i=e_k(m_i)$. Es determinista: bloques de texto plano iguales producen cifrados iguales
  $\Rightarrow$ \textbf{revela patrones/estructura} y \textbf{no es CPA-Secure}. No debe
  usarse.
  \item \textbf{HTTPS = HTTP sobre TLS.} El TLS Handshake acuerda un método de intercambio de
  clave (RSA, Diffie--Hellman con certificado, ECDH, etc.), autentica con firmas y deriva un
  \emph{Master Secret}; la criptografía negociable para el resto de la sesión incluye cifrado
  simétrico (recomendado AES-GCM) y HMAC/AEAD para integridad.
  \item \textbf{Diffie--Hellman} (Clase 5): acuerdo de clave cuya seguridad descansa en el
  problema del logaritmo discreto (DLP); permite que ambas partes compartan una clave secreta
  sobre un canal inseguro sin haberla intercambiado antes. \textbf{DH puro no autentica}: es
  vulnerable a Man-in-the-Middle si no se combina con certificados (por eso en TLS va con
  RSA/DSS).
  \item \textbf{Cifrado autenticado / AEAD} (Clase 3): en la práctica, si se necesita
  confidencialidad casi siempre también se necesita integridad, porque el cifrado simple
  (como ECB) puede ser manipulado sin que nadie lo note. La recomendación es usar un modo
  \textbf{AEAD} como \textbf{AES-GCM} (o CCM), que da confidencialidad e integridad juntas.
\end{itemize}

\paragraph{Notas y conexiones.}
El escenario combina dos fallas de diseño típicas de la materia: (1) enviar una clave en
claro por un canal no cifrado (HTTP) para luego usarla en un cifrado simétrico, y (2) usar
ECB, el modo que la Clase 2 marca como ``NO utilizar'' por no ser CPA-Secure y no dar
integridad. Conecta con la distinción Clase 3 entre cifrado (confidencialidad) y MAC/AEAD
(integridad): un cifrado por sí solo, aunque fuera un modo seguro como CBC o GCM sin
autenticar, no impide que un atacante manipule bits del cifrado sin ser detectado si no hay
un mecanismo de integridad de por medio.

\paragraph{Resolución.}

\textbf{a) No, no necesariamente se hubieran solucionado todos los problemas.}
\begin{itemize}
  \item Usando \textbf{HTTPS con TLS} para el paso donde la central envía la llave, sí se
  soluciona el problema de la \textbf{comunicación de la llave}: TLS resuelve el intercambio
  de claves de forma autenticada (certificados) y cifra el canal, evitando que un atacante
  intercepte la llave en texto plano durante el handshake inicial. Esto ataca directamente la
  falla ``fallas en la comunicación de la llave''.
  \item Sin embargo, HTTPS solo protege la \textbf{comunicación} del dispositivo con la
  central (la llave viajando por HTTP). \textbf{No cambia el método de encripción que el
  dispositivo usa para enviar lo recolectado}, que sigue siendo \textbf{AES-ECB}. Como ECB es
  determinista y no es CPA-Secure, \textbf{no da integridad} (bloques de datos repetidos
  producen cifrados idénticos, y un atacante puede reordenar/reemplazar bloques de cifrado
  sin que el receptor lo detecte, ya que ECB no incluye ningún mecanismo de autenticación).
  Por lo tanto el problema de \textbf{integridad en la comunicación de los dispositivos}
  \textbf{no queda resuelto} solo con HTTPS en el paso de intercambio de la llave.
\end{itemize}

\textbf{b) Diffie--Hellman.} Si no se puede pagar el costo de HTTPS/TLS completo, se puede
usar directamente un protocolo de \textbf{intercambio de clave Diffie--Hellman}: permite que
la central y el dispositivo acuerden una clave secreta compartida sobre un canal inseguro sin
haberla compartido antes, basando su seguridad en el problema del logaritmo discreto. Queda
como salvedad (marcada en el material) que \textbf{DH puro no autentica a las partes} y es
susceptible a un ataque Man-in-the-Middle si un atacante se interpone en el intercambio; para
mitigarlo haría falta combinarlo con algún mecanismo de autenticación (certificados o claves
públicas ya conocidas), tal como TLS lo hace con RSA/DSS.

\textbf{c)} Para asegurar la integridad en el envío de la información recolectada hacia la
central hace falta:
\begin{enumerate}
  \item \textbf{Dejar de usar ECB en solitario:} es un modo sin ningún mecanismo de
  autenticación, determinista y no CPA-Secure.
  \item \textbf{Pasar a un esquema de cifrado autenticado (AEAD)}, que da confidencialidad
  \emph{e} integridad en una sola construcción: \textbf{AES-GCM} (recomendación práctica del
  material: rápido, estandarizado, en todas las librerías serias) o \textbf{AES-CCM}. Ambos
  agregan, además del cifrado, un \emph{tag} de autenticación que permite al receptor
  verificar que el mensaje cifrado no fue modificado en el camino.
  \item Alternativamente (o en conjunto), agregar explícitamente un \textbf{MAC/HMAC} sobre lo
  enviado, de forma que la central pueda \textbf{chequear que el mensaje no fue modificado en
  el camino} (verificar el tag con la clave compartida antes de aceptar la medición).
\end{enumerate}
Nota de fidelidad: el material no detalla un protocolo específico de bajo costo distinto de
Diffie--Hellman para IoT (p.\ ej. protocolos livianos de sensores); la respuesta se limita a
lo cubierto en las clases de Protocolos y MACs/Cifrado Autenticado.

%=====================================================================
\subsection{Ejercicio 2 --- Ocultar un mensaje ``tipo sombras'' (OTP + integridad con hash/HMAC)}
%=====================================================================

\begin{consigna}
Se quiere ocultar un mensaje de forma segura con una técnica similar a las sombras. Dado un
mensaje ``m'' se eligen dos números random ($x_1,x_2$) y se generan por dos canales separados
(posiblemente en lugares distintos):
\begin{itemize}
  \item $p_1 = x_1 \oplus m \oplus x_2$
  \item $p_2 = x_2 \oplus x_1$
  \item $p_1$ y $h(p_2)$
  \item $p_2$ y $h(p_1)$
\end{itemize}
\begin{enumerate}
  \item[a)] Explicar por qué si un atacante tiene acceso a solo 1 canal de comunicación aun se
  mantiene la confidencialidad.
  \item[b)] Explicar por qué si un atacante tiene acceso a un solo canal de comunicación aun se
  mantiene la integridad.
  \item[c)] Si un atacante tiene acceso a los 2 canales de información, cómo hacer para que se
  mantenga \underline{solo la integridad}.
\end{enumerate}
\end{consigna}

\paragraph{Temas.}
Primera Parte, Clase 2 --- Cifrado (\S secreto perfecto, \S One-Time Pad / Vernam); Clase 5 ---
Protocolos (\S secreto compartido de Shamir, umbral $(t,n)$: con menos del umbral no se filtra
información); Clase 3 --- MACs y Cifrado Autenticado (\S funciones de hash, \S HMAC,
\S no-repudio: por qué MAC/hash no lo dan).

\paragraph{Qué necesitás saber.}
\begin{itemize}
  \item \textbf{OTP (Clase 2):} $e_k(m) = m \oplus k$ con $k$ uniforme y del mismo tamaño que
  $m$ tiene \textbf{secreto perfecto}: $\Pr[M{=}m\mid C{=}c] = \Pr[M{=}m]$. Si $x_1$ y $x_2$
  son aleatorios e independientes, $x_1\oplus x_2$ se comporta como una clave uniforme.
  \item \textbf{Secreto compartido / esquema umbral (Clase 5):} repartir información en
  ``sombras'' tal que \textbf{con menos de las necesarias no se filtra ninguna información}
  del secreto (entropía máxima). Acá el mensaje se reparte en dos partes ($p_1,p_2$) tal que
  ninguna por separado revela $m$.
  \item \textbf{Función de hash (Clase 3):} \textbf{no tiene clave} (es información pública);
  depende de todos los bits del mensaje (efecto avalancha) y permite detectar si el mensaje
  fue alterado comparando el hash recibido contra el hash recalculado. Al no tener clave, un
  atacante que \textbf{controla el mensaje y puede recalcular el hash él mismo} puede
  modificar ambos de forma consistente sin ser detectado.
  \item \textbf{HMAC (Clase 3):} MAC construido a partir de una función de hash y una
  \textbf{clave} $k$ compartida; a diferencia del hash simple, el atacante sin la clave
  \textbf{no puede recalcular un tag válido} aunque conozca y modifique el mensaje.
\end{itemize}

\paragraph{Notas y conexiones.}
El esquema es análogo a un \textbf{secreto compartido de 2 sombras} (Clase 5) combinado con
\textbf{OTP} (Clase 2): $p_1$ es el mensaje enmascarado con el XOR de dos valores aleatorios
(equivalente a una clave OTP $x_1\oplus x_2$), y $p_2$ es esa misma máscara sin el mensaje.
Ninguna de las dos partes por sí sola alcanza para reconstruir $m$: hace falta juntar ambos
canales, igual que con el umbral de Shamir se necesita juntar $t$ sombras. El agregado de
$h(\cdot)$ cruzado conecta con Clase 3 (funciones de hash sin clave) y anticipa la necesidad
de HMAC para integridad \emph{autenticada} cuando el atacante controla todo.

\paragraph{Resolución.}

\textbf{a) Confidencialidad con acceso a un solo canal.}
\begin{itemize}
  \item \textbf{Canal 1} entrega $(p_1, h(p_2))$, con $p_1 = x_1\oplus m\oplus x_2$. Como
  $x_1$ y $x_2$ son elegidos al azar e independientes, $k = x_1\oplus x_2$ se comporta como
  una \textbf{clave OTP uniforme}: $p_1 = m \oplus k$. Por el resultado de secreto perfecto
  del OTP, sin conocer $k$ el atacante no obtiene ninguna información sobre $m$ a partir de
  $p_1$ solo. $h(p_2)$ tampoco aporta nada sobre $m$ (depende solo de $p_2 = x_1\oplus x_2$,
  que no involucra a $m$).
  \item \textbf{Canal 2} entrega $(p_2, h(p_1))$, con $p_2 = x_1\oplus x_2$. Esto es
  exactamente la máscara $k$ usada para enmascarar $m$, pero \textbf{sin el mensaje}: no
  contiene ninguna información de $m$. $h(p_1)$ tampoco revela $m$ porque una función de hash
  no es invertible (no permite recuperar el mensaje original a partir del resumen).
\end{itemize}
En ambos casos se necesitan \textbf{ambos canales} para calcular
$p_1 \oplus p_2 = (x_1\oplus m\oplus x_2)\oplus(x_1\oplus x_2) = m$; con uno solo la
confidencialidad de $m$ se mantiene.

\textbf{b) Integridad con acceso a un solo canal.}
Cada canal lleva el valor a transmitir junto con el \textbf{hash del otro} valor
($p_1$ con $h(p_2)$, y $p_2$ con $h(p_1)$), es decir, hashes \textbf{cruzados}. Si el
atacante solo controla un canal (por ejemplo el que lleva $p_1$ y $h(p_2)$) y modifica $p_1$
en tránsito, no tiene forma de recalcular un $h(p_2)$ consistente porque \textbf{no tiene
acceso a $p_2$} (viaja por el otro canal, que no controla). Al llegar los dos valores al
receptor, este recalcula $h(p_2)$ recibido en el otro canal y $h(p_1)$ recibido en este, y
si $p_1$ fue alterado, el hash recalculado de $p_1$ no coincidirá con el $h(p_1)$ que viajó
por el canal 2, delatando la alteración. Con un solo canal comprometido, la verificación
cruzada contra el otro canal (no comprometido) detecta la manipulación.

\textbf{c) Con acceso a los dos canales, mantener solo la integridad.}
Si el atacante controla ambos canales, puede leer $p_1$ y $p_2$ enteros, calcular
$h(\cdot)$ él mismo (la función de hash \textbf{no usa clave}, es pública) y modificar
$p_1$ y/o $p_2$ recalculando los hashes cruzados de forma consistente: el esquema de hash
simple deja de servir para integridad en ese caso (y de hecho la confidencialidad de $m$
también se pierde, porque con los dos canales se recupera $m = p_1\oplus p_2$; por eso la
consigna pide mantener \emph{solo} la integridad, asumiendo que la confidencialidad ya no es
sostenible con acceso total a ambos canales).

Para mantener la integridad hay que reemplazar el hash simple $h(\cdot)$ por un mecanismo
\textbf{que dependa de una clave que el atacante no tenga}:
\begin{itemize}
  \item Usar \textbf{HMAC} en lugar de un hash sin clave: $\mathrm{HMAC}_k(p_2)$ en el canal 1
  y $\mathrm{HMAC}_k(p_1)$ en el canal 2, con una clave $k$ secreta compartida entre emisor y
  receptor mediante un canal previo seguro. Sin conocer $k$, el atacante no puede generar un
  tag $\mathrm{HMAC}_k$ válido para un $p_1$ o $p_2$ modificado, aunque vea y edite el
  contenido de ambos canales.
  \item Alternativamente, \textbf{firmar digitalmente} $p_1$ y $p_2$ con la clave
  \textbf{privada} del emisor: el atacante, aun controlando ambos canales, no puede falsificar
  la firma sin esa clave privada, por lo que cualquier alteración se detecta al verificar la
  firma con la clave pública correspondiente. (Esta variante, a diferencia de HMAC, además
  daría no-repudio, ya que el material remarca que un MAC/HMAC simétrico \textbf{no} da
  no-repudio porque ambas partes comparten la misma clave.)
\end{itemize}

%=====================================================================
\subsection{Ejercicio 3 --- Pentesting a una página de noticias (hipótesis de falla)}
%=====================================================================

\begin{consigna}
Una página de noticias en la que hay un usuario y un generador de noticias. A cada noticia se
le puede poner etiquetas como categorías (Deportes, Interés general) o Ubicación (Argentina).

Cada usuario lleva un puntaje (+1 por cada noticia rateada positivamente y -1 por cada noticia
rateada negativamente). La página tiene una API REST y los llamados a la misma se hacen por
HTTPS.

Se le pagará por visualizaciones al generador de la noticia teniendo en cuenta algunos
factores que no me acuerdo bien y no venían al caso.

Se hizo un análisis de la página y se encontró que:
\begin{enumerate}
  \item Existe una URL para pedir todas las noticias:
  \texttt{https://dominio.com/news?limit=100\&page=20}
  \item Cuando se le recomienda o no recomienda una noticia se hace a través de
  \texttt{https://dominio.com/\{generator\_id\}/\{news\_id\}/upvoted} o
  \texttt{https://dominio.com/\{generator\_id\}/\{news\_id\}/downvoted}
  \item Las personas se registran por usuario y password
  \item Se chequearon los inputs de la página y cualquier intento de tag HTML ingresado
  aparecerá como texto en la página $\rightarrow$ XSS no probable.
\end{enumerate}
Establecer 4 hipótesis de falla con las respectivas pruebas que haría.
\end{consigna}

\paragraph{Temas.}
Segunda Parte, Clase 10b --- Pentesting (\S Metodología de Hipótesis de Falla: los pasos en
orden --- Recolección, Hipótesis, Prueba, Generalización, Eliminación; \S Injection flaws;
\S OWASP); Clase 8 --- Principios de Diseño y Vulnerabilidades (relacionado: validación de
entrada, XSS/CWE-79).

\paragraph{Qué necesitás saber.}
\begin{itemize}
  \item La \textbf{metodología de hipótesis de falla} (Clase 10b) tiene 5 pasos en orden:
  \textbf{Recolección de información} $\to$ \textbf{Planteo de hipótesis de falla} $\to$
  \textbf{Prueba} $\to$ \textbf{Generalización} $\to$ \textbf{Eliminación} (opcional). Cada
  hipótesis debe: (1) partir de la información recolectada del sistema (los 4 hallazgos del
  enunciado), (2) suponer una vulnerabilidad concreta, y (3) tener una \textbf{prueba}
  asociada que la verifique.
  \item El material define el \textbf{XSS (CWE-79)} como inyección de JS ejecutado por el
  navegador de la víctima, y aquí el enunciado ya lo descarta explícitamente (los tags HTML
  se muestran como texto, no se ejecutan) --- es un hallazgo del paso de Recolección que cierra
  esa hipótesis.
  \item El material menciona en general \textbf{OWASP} como organización que publica el
  Top 10 de vulnerabilidades web y menciona \textbf{Injection flaws} (SQLi, XSS, OS command)
  como la familia más frecuente ($\sim$90\% según el profesor).
\end{itemize}

\begin{warnbox}{Aviso de fidelidad al material}
El enunciado del alumno (usado solo como referencia) etiqueta las hipótesis con nombres de
categorías estilo OWASP Top 10 API (``Improper Authentication'', ``Broken Access Control'',
``BOLA'', ataques de \emph{replay}). Los resúmenes de la materia \textbf{sí} cubren OWASP
como organización/Top 10 y la familia de \emph{Injection flaws}, y \textbf{sí} cubren
conceptos de control de acceso/autenticación de otras clases, pero \textbf{no} desarrollan
esas siglas puntuales (BOLA, Improper Authentication, Broken Access Control como categorías
nombradas) como parte del temario de Pentesting. Por eso, en la resolución de abajo se
plantean las 4 hipótesis con el vocabulario y la metodología efectivamente vistos en la
Clase 10b (hipótesis de falla + prueba), evitando nombrar categorías OWASP específicas que no
están en el material, y describiendo la falla en términos genéricos de control de acceso /
autenticación / lógica de negocio ya cubiertos en Segunda Parte.
\end{warnbox}

\paragraph{Resolución.}
Aplicando la metodología de hipótesis de falla (Clase 10b) sobre los 4 hallazgos de la
Recolección de información:

\begin{enumerate}
  \item \textbf{Hipótesis 1 --- falta de límite de intentos en el login.} El sistema se
  registra y autentica por usuario y contraseña (hallazgo 3), y no hay ninguna mención de
  control sobre los intentos de autenticación. \textbf{Hipótesis:} el endpoint de login no
  limita ni registra intentos fallidos, permitiendo probar contraseñas por fuerza bruta contra
  una cuenta válida. \textbf{Prueba:} realizar múltiples intentos de autenticación contra un
  usuario existente con contraseñas distintas y observar si el sistema bloquea, demora o
  registra los intentos repetidos, o si los permite indefinidamente.

  \item \textbf{Hipótesis 2 --- falta de verificación del vínculo usuario/\texttt{generator\_id}
  al votar.} Los endpoints de recomendación usan \texttt{\{generator\_id\}/\{news\_id\}/upvoted}
  y \texttt{downvoted} directamente en la URL (hallazgo 2), sin mostrar cómo se valida contra
  quién está autenticado. \textbf{Hipótesis:} el back-end no verifica que el
  \texttt{generator\_id} de la URL corresponda al usuario autenticado o autorizado para votar,
  permitiendo que un usuario vote (sume/reste puntaje) en nombre de otro generador o sin estar
  autenticado. \textbf{Prueba:} (i) sin loguearse, llamar directamente a la URL de
  \texttt{upvoted}/\texttt{downvoted} con un \texttt{generator\_id}/\texttt{news\_id} válidos y
  observar si el sistema restringe el acceso; (ii) logueado como un usuario, llamar a la misma
  URL con un \texttt{generator\_id} que no es el propio y observar si el sistema lo permite.

  \item \textbf{Hipótesis 3 --- manipulación del parámetro \texttt{limit} de la URL de
  noticias.} El hallazgo 1 expone una URL con parámetros \texttt{limit} y \texttt{page}
  controlados por el cliente. \textbf{Hipótesis:} el valor de \texttt{limit} no está validado
  ni acotado en el servidor, permitiendo pedir una cantidad arbitrariamente grande de noticias
  en una sola request (degradación del servicio / consumo excesivo de recursos).
  \textbf{Prueba:} modificar el valor de \texttt{limit} a números extremadamente altos (o
  negativos/no numéricos) y observar la respuesta y el tiempo/recursos que consume el
  servidor.

  \item \textbf{Hipótesis 4 --- falta de control sobre votos repetidos del mismo usuario.} El
  puntaje del generador depende de +1/-1 por cada rating (hallazgo del enunciado principal).
  \textbf{Hipótesis:} el sistema no valida que un mismo usuario ya haya votado una misma
  noticia, permitiendo repetir el llamado a \texttt{upvoted}/\texttt{downvoted} para inflar o
  deflar artificialmente el puntaje de un generador (equivalente a un ataque de repetición
  sobre la misma acción). \textbf{Prueba:} repetir la misma llamada de
  \texttt{upvoted}/\texttt{downvoted} (mismo usuario, misma noticia) dos o más veces y verificar
  si el puntaje del generador se incrementa cada vez o si el sistema detecta y descarta el
  duplicado.
\end{enumerate}

Nota: el hallazgo 4 del enunciado (chequeo de tags HTML) queda \textbf{descartado} como
hipótesis de XSS por el propio enunciado (paso de Recolección ya cerrado): cualquier tag
ingresado se muestra como texto, no se ejecuta, así que no aplica la familia de
\emph{Injection flaws} tipo XSS sobre ese input en particular.
